\chapter{Persistent homology in plain language}
\label{ch:ph}
\label{ch:persistent-homology}

This chapter is the mathematical core of the document.
We build from spatial intuitions that any linguist already has---the sense
that some things cluster, some things form loops, some structures have
holes---and turn those intuitions into machinery precise enough to run
on attention graphs.
The tools come from algebraic topology, but the ideas are
concrete: count the pieces, count the holes, and track how both change
as we zoom in.

\bigskip

\begin{table}[ht]
\centering
\caption{Core notation used throughout this chapter and the rest of the document.}
\label{tab:notation}
\begin{tabular}{c p{10cm}}
\hline
\textbf{Symbol} & \textbf{Meaning} \\
\hline
$\varepsilon$ & Filtration parameter (threshold). As $\varepsilon$ grows, more simplices enter the complex. \\
$W$ & Attention weight matrix from a single head. $W_{ij}$ is the attention weight from token $i$ to token $j$. \\
$d(u,v)$ & Distance between two points (or tokens). In our pipeline, typically derived from attention weights. \\
$H_0$ & Zeroth homology group: its rank counts connected components. \\
$H_1$ & First homology group: its rank counts independent loops (1-cycles). \\
$H_2$ & Second homology group: its rank counts enclosed voids. Rarely informative for attention graphs. \\
$\beta_0, \beta_1$ & Betti numbers: $\beta_k = \mathrm{rank}(H_k)$. $\beta_0$ is the number of components; $\beta_1$ is the number of loops. \\
$(b_i, d_i)$ & A persistence pair: feature $i$ is born at filtration value $b_i$ and dies at $d_i$. \\
$d_B$ & Bottleneck distance between two persistence diagrams. \\
$W_p$ & $p$-Wasserstein distance between two persistence diagrams. \\
\hline
\end{tabular}
\end{table}

\section{The shape of data}
\label{sec:shape-of-data}

Suppose we have a set of points in some space and we want to describe
their \emph{shape}---not the coordinates of individual points, but the
large-scale geometric structure they form. Do they fall into separate
clusters? Do they trace out a ring? Is there a void in the middle?

These are topological questions.
Topology studies the properties of a space that survive continuous
deformation: stretching and bending are allowed, tearing and gluing are
not.
A coffee mug and a donut are topologically identical (both have exactly
one hole); a coffee mug and a bowling ball are not.
\textcite{carlsson2009} argued that the same invariants that
distinguish the mug from the ball can distinguish genuinely different
shapes in high-dimensional data, and that this perspective---topology
applied to finite, noisy point clouds---constitutes a new kind of data
analysis.

The field that grew from that argument is \emph{topological data
analysis} (TDA), and its central tool is \emph{persistent homology}
\parencite{edelsbrunner2010, zomorodian2005}.
We use persistent homology in this project because we care about
\emph{structural} differences between attention graphs, not just
differences in individual edge weights.
Two attention heads could have very different weight distributions yet
produce the same graph shape (same clusters, same loops); two other
heads could have similar weight distributions yet produce topologically
distinct graphs.
Persistent homology sees the distinction that raw statistics miss.


\section{Simplicial complexes: the building blocks}
\label{sec:simplicial-complexes}

To do topology on a finite point set, we need to build a space from it.
The standard construction uses \emph{simplicial complexes}---objects
assembled from a small kit of geometric pieces:

\begin{itemize}
  \item A \textbf{0-simplex} is a point (a vertex).
  \item A \textbf{1-simplex} is an edge connecting two vertices.
  \item A \textbf{2-simplex} is a filled triangle on three vertices.
  \item In general, a \textbf{$k$-simplex} is the convex hull of $k+1$
        points---but for attention graphs we rarely go above $k = 2$.
\end{itemize}

A \emph{simplicial complex} $K$ is a collection of simplices that is
closed under taking faces: if a triangle belongs to $K$, then so do its
three edges and three vertices.
The complex is the ``space'' we do topology on.

How do we get a simplicial complex from data?
Given a set of points with pairwise distances, we pick a threshold
$\varepsilon$ and connect every pair of points whose distance is at most
$\varepsilon$.
This gives us a graph (vertices and edges).
Whenever three points are mutually connected, we fill in the triangle.
Whenever four points are mutually connected, we fill in the tetrahedron.
And so on.
The result is the \emph{Vietoris--Rips complex} at scale $\varepsilon$,
written $\mathrm{VR}_\varepsilon$.

For attention graphs, the ``points'' are tokens and the ``distances''
come from the attention weight matrix $W$ (specifically, from a
transformation of the weights---\autoref{ch:pipeline} details the
conversion).
At a low threshold only the strongest connections appear;
at a high threshold nearly everything is connected.
The interesting structure lives in between.


\section{Filtrations: watching structure emerge}
\label{sec:filtration}

A single snapshot of the Vietoris--Rips complex at one threshold
$\varepsilon$ captures some of the data's shape, but choosing the
``right'' $\varepsilon$ is both arbitrary and lossy.
Persistent homology sidesteps the problem by refusing to choose.
Instead, we sweep $\varepsilon$ from $0$ to $\infty$ and watch the
complex grow.

A \emph{filtration} is this nested sequence of complexes:
\[
  \emptyset
  = K_0 \subseteq K_1 \subseteq K_2 \subseteq \cdots \subseteq K_m = K,
\]
where each $K_i$ is the Vietoris--Rips complex at some threshold
$\varepsilon_i$, and $\varepsilon_0 < \varepsilon_1 < \cdots <
\varepsilon_m$.
As $\varepsilon$ increases, simplices appear and never disappear---once
an edge is added, it stays.

% filtration.tex — Vietoris-Rips filtration on a small point cloud
% Include via % filtration.tex — Vietoris-Rips filtration on a small point cloud
% Include via % filtration.tex — Vietoris-Rips filtration on a small point cloud
% Include via \input{figures/filtration} from overview.tex

\begin{figure}[H]
\centering
\begin{tikzpicture}[
    point/.style={circle, fill=black, inner sep=1.8pt},
    edge/.style={thick, blue!60!black},
    tri/.style={fill=blue!15, draw=none},
    panel label/.style={font=\small\bfseries, anchor=north},
    eps label/.style={font=\small, anchor=north},
    >=Stealth
]

% --- Point coordinates (shared across all four panels) ---
% Six points arranged in a rough cluster with one outlier
\newcommand{\drawpoints}[1]{
    \node[point] (A#1) at ($(#1,0) + (0.0, 0.8)$) {};
    \node[point] (B#1) at ($(#1,0) + (0.7, 1.5)$) {};
    \node[point] (C#1) at ($(#1,0) + (1.4, 0.9)$) {};
    \node[point] (D#1) at ($(#1,0) + (0.3, 0.0)$) {};
    \node[point] (E#1) at ($(#1,0) + (1.1, 0.1)$) {};
    \node[point] (F#1) at ($(#1,0) + (2.6, 1.2)$) {};
}

% ---- Panel 1: points only, varepsilon = 0 ----
\drawpoints{0}
\node[panel label] at (0.8, -0.5) {(a)};
\node[eps label] at (0.8, -0.9) {$\varepsilon = 0$};

% ---- Panel 2: a few short edges, varepsilon = 0.9 ----
\drawpoints{4}
\draw[edge] (A4) -- (D4);
\draw[edge] (D4) -- (E4);
\draw[edge] (A4) -- (B4);
\node[panel label] at (4.8, -0.5) {(b)};
\node[eps label] at (4.8, -0.9) {$\varepsilon = 0.9$};

% ---- Panel 3: more edges + a filled triangle, varepsilon = 1.2 ----
\drawpoints{8}
% filled triangle A-D-E
\fill[tri] (A8.center) -- (D8.center) -- (E8.center) -- cycle;
\draw[edge] (A8) -- (D8);
\draw[edge] (D8) -- (E8);
\draw[edge] (A8) -- (E8);
\draw[edge] (A8) -- (B8);
\draw[edge] (A8) -- (C8);
\draw[edge] (B8) -- (C8);
\draw[edge] (C8) -- (E8);
\node[panel label] at (8.8, -0.5) {(c)};
\node[eps label] at (8.8, -0.9) {$\varepsilon = 1.2$};

% ---- Panel 4: all edges, varepsilon = 2.0 ----
\drawpoints{12}
% filled triangles
\fill[tri] (A12.center) -- (D12.center) -- (E12.center) -- cycle;
\fill[tri] (A12.center) -- (B12.center) -- (C12.center) -- cycle;
\fill[tri] (A12.center) -- (C12.center) -- (E12.center) -- cycle;
\fill[tri] (C12.center) -- (E12.center) -- (F12.center) -- cycle;
% all edges
\draw[edge] (A12) -- (B12);
\draw[edge] (A12) -- (C12);
\draw[edge] (A12) -- (D12);
\draw[edge] (A12) -- (E12);
\draw[edge] (B12) -- (C12);
\draw[edge] (B12) -- (F12);
\draw[edge] (C12) -- (E12);
\draw[edge] (C12) -- (F12);
\draw[edge] (D12) -- (E12);
\draw[edge] (E12) -- (F12);
\node[panel label] at (12.8, -0.5) {(d)};
\node[eps label] at (12.8, -0.9) {$\varepsilon = 2.0$};

% ---- Arrow showing direction ----
\draw[->, thick, gray] (2.2, -1.4) -- (13.8, -1.4)
    node[midway, below, font=\small, gray] {$\varepsilon$ increases};

\end{tikzpicture}
\caption{%
    Vietoris-Rips filtration on six points at increasing thresholds
    $\varepsilon$. At $\varepsilon = 0$ only the points exist. As
    $\varepsilon$ grows, edges appear between points whose distance is at
    most $\varepsilon$, and filled triangles (2-simplices) appear when all
    three pairwise edges are present. Tracking how connected components
    merge ($H_0$) and loops form and fill ($H_1$) across the filtration is
    the core idea of persistent homology.%
}
\label{fig:filtration}
\end{figure}
 from overview.tex

\begin{figure}[H]
\centering
\begin{tikzpicture}[
    point/.style={circle, fill=black, inner sep=1.8pt},
    edge/.style={thick, blue!60!black},
    tri/.style={fill=blue!15, draw=none},
    panel label/.style={font=\small\bfseries, anchor=north},
    eps label/.style={font=\small, anchor=north},
    >=Stealth
]

% --- Point coordinates (shared across all four panels) ---
% Six points arranged in a rough cluster with one outlier
\newcommand{\drawpoints}[1]{
    \node[point] (A#1) at ($(#1,0) + (0.0, 0.8)$) {};
    \node[point] (B#1) at ($(#1,0) + (0.7, 1.5)$) {};
    \node[point] (C#1) at ($(#1,0) + (1.4, 0.9)$) {};
    \node[point] (D#1) at ($(#1,0) + (0.3, 0.0)$) {};
    \node[point] (E#1) at ($(#1,0) + (1.1, 0.1)$) {};
    \node[point] (F#1) at ($(#1,0) + (2.6, 1.2)$) {};
}

% ---- Panel 1: points only, varepsilon = 0 ----
\drawpoints{0}
\node[panel label] at (0.8, -0.5) {(a)};
\node[eps label] at (0.8, -0.9) {$\varepsilon = 0$};

% ---- Panel 2: a few short edges, varepsilon = 0.9 ----
\drawpoints{4}
\draw[edge] (A4) -- (D4);
\draw[edge] (D4) -- (E4);
\draw[edge] (A4) -- (B4);
\node[panel label] at (4.8, -0.5) {(b)};
\node[eps label] at (4.8, -0.9) {$\varepsilon = 0.9$};

% ---- Panel 3: more edges + a filled triangle, varepsilon = 1.2 ----
\drawpoints{8}
% filled triangle A-D-E
\fill[tri] (A8.center) -- (D8.center) -- (E8.center) -- cycle;
\draw[edge] (A8) -- (D8);
\draw[edge] (D8) -- (E8);
\draw[edge] (A8) -- (E8);
\draw[edge] (A8) -- (B8);
\draw[edge] (A8) -- (C8);
\draw[edge] (B8) -- (C8);
\draw[edge] (C8) -- (E8);
\node[panel label] at (8.8, -0.5) {(c)};
\node[eps label] at (8.8, -0.9) {$\varepsilon = 1.2$};

% ---- Panel 4: all edges, varepsilon = 2.0 ----
\drawpoints{12}
% filled triangles
\fill[tri] (A12.center) -- (D12.center) -- (E12.center) -- cycle;
\fill[tri] (A12.center) -- (B12.center) -- (C12.center) -- cycle;
\fill[tri] (A12.center) -- (C12.center) -- (E12.center) -- cycle;
\fill[tri] (C12.center) -- (E12.center) -- (F12.center) -- cycle;
% all edges
\draw[edge] (A12) -- (B12);
\draw[edge] (A12) -- (C12);
\draw[edge] (A12) -- (D12);
\draw[edge] (A12) -- (E12);
\draw[edge] (B12) -- (C12);
\draw[edge] (B12) -- (F12);
\draw[edge] (C12) -- (E12);
\draw[edge] (C12) -- (F12);
\draw[edge] (D12) -- (E12);
\draw[edge] (E12) -- (F12);
\node[panel label] at (12.8, -0.5) {(d)};
\node[eps label] at (12.8, -0.9) {$\varepsilon = 2.0$};

% ---- Arrow showing direction ----
\draw[->, thick, gray] (2.2, -1.4) -- (13.8, -1.4)
    node[midway, below, font=\small, gray] {$\varepsilon$ increases};

\end{tikzpicture}
\caption{%
    Vietoris-Rips filtration on six points at increasing thresholds
    $\varepsilon$. At $\varepsilon = 0$ only the points exist. As
    $\varepsilon$ grows, edges appear between points whose distance is at
    most $\varepsilon$, and filled triangles (2-simplices) appear when all
    three pairwise edges are present. Tracking how connected components
    merge ($H_0$) and loops form and fill ($H_1$) across the filtration is
    the core idea of persistent homology.%
}
\label{fig:filtration}
\end{figure}
 from overview.tex

\begin{figure}[H]
\centering
\begin{tikzpicture}[
    point/.style={circle, fill=black, inner sep=1.8pt},
    edge/.style={thick, blue!60!black},
    tri/.style={fill=blue!15, draw=none},
    panel label/.style={font=\small\bfseries, anchor=north},
    eps label/.style={font=\small, anchor=north},
    >=Stealth
]

% --- Point coordinates (shared across all four panels) ---
% Six points arranged in a rough cluster with one outlier
\newcommand{\drawpoints}[1]{
    \node[point] (A#1) at ($(#1,0) + (0.0, 0.8)$) {};
    \node[point] (B#1) at ($(#1,0) + (0.7, 1.5)$) {};
    \node[point] (C#1) at ($(#1,0) + (1.4, 0.9)$) {};
    \node[point] (D#1) at ($(#1,0) + (0.3, 0.0)$) {};
    \node[point] (E#1) at ($(#1,0) + (1.1, 0.1)$) {};
    \node[point] (F#1) at ($(#1,0) + (2.6, 1.2)$) {};
}

% ---- Panel 1: points only, varepsilon = 0 ----
\drawpoints{0}
\node[panel label] at (0.8, -0.5) {(a)};
\node[eps label] at (0.8, -0.9) {$\varepsilon = 0$};

% ---- Panel 2: a few short edges, varepsilon = 0.9 ----
\drawpoints{4}
\draw[edge] (A4) -- (D4);
\draw[edge] (D4) -- (E4);
\draw[edge] (A4) -- (B4);
\node[panel label] at (4.8, -0.5) {(b)};
\node[eps label] at (4.8, -0.9) {$\varepsilon = 0.9$};

% ---- Panel 3: more edges + a filled triangle, varepsilon = 1.2 ----
\drawpoints{8}
% filled triangle A-D-E
\fill[tri] (A8.center) -- (D8.center) -- (E8.center) -- cycle;
\draw[edge] (A8) -- (D8);
\draw[edge] (D8) -- (E8);
\draw[edge] (A8) -- (E8);
\draw[edge] (A8) -- (B8);
\draw[edge] (A8) -- (C8);
\draw[edge] (B8) -- (C8);
\draw[edge] (C8) -- (E8);
\node[panel label] at (8.8, -0.5) {(c)};
\node[eps label] at (8.8, -0.9) {$\varepsilon = 1.2$};

% ---- Panel 4: all edges, varepsilon = 2.0 ----
\drawpoints{12}
% filled triangles
\fill[tri] (A12.center) -- (D12.center) -- (E12.center) -- cycle;
\fill[tri] (A12.center) -- (B12.center) -- (C12.center) -- cycle;
\fill[tri] (A12.center) -- (C12.center) -- (E12.center) -- cycle;
\fill[tri] (C12.center) -- (E12.center) -- (F12.center) -- cycle;
% all edges
\draw[edge] (A12) -- (B12);
\draw[edge] (A12) -- (C12);
\draw[edge] (A12) -- (D12);
\draw[edge] (A12) -- (E12);
\draw[edge] (B12) -- (C12);
\draw[edge] (B12) -- (F12);
\draw[edge] (C12) -- (E12);
\draw[edge] (C12) -- (F12);
\draw[edge] (D12) -- (E12);
\draw[edge] (E12) -- (F12);
\node[panel label] at (12.8, -0.5) {(d)};
\node[eps label] at (12.8, -0.9) {$\varepsilon = 2.0$};

% ---- Arrow showing direction ----
\draw[->, thick, gray] (2.2, -1.4) -- (13.8, -1.4)
    node[midway, below, font=\small, gray] {$\varepsilon$ increases};

\end{tikzpicture}
\caption{%
    Vietoris-Rips filtration on six points at increasing thresholds
    $\varepsilon$. At $\varepsilon = 0$ only the points exist. As
    $\varepsilon$ grows, edges appear between points whose distance is at
    most $\varepsilon$, and filled triangles (2-simplices) appear when all
    three pairwise edges are present. Tracking how connected components
    merge ($H_0$) and loops form and fill ($H_1$) across the filtration is
    the core idea of persistent homology.%
}
\label{fig:filtration}
\end{figure}


\autoref{fig:filtration} illustrates this on a small example.
At the smallest thresholds we see only isolated vertices (many connected
components, no loops).
As $\varepsilon$ grows, edges appear and components merge.
At some point a loop may form: three or more edges close into a cycle
without a filled triangle inside.
If $\varepsilon$ grows further and the triangle fills in, the loop
vanishes.
The filtration records the entire history of these births and deaths.


\section{Homology: counting the holes}
\label{sec:homology}

Homology is the algebraic machinery that counts topological
features---components, loops, voids---of a simplicial complex.
We will not develop the full algebraic theory (which requires chain
complexes, boundary operators, and quotient groups; see
\textcite{edelsbrunner2010} for the complete treatment).
What matters for us is the output: for each dimension $k$, the
\emph{$k$-th homology group} $H_k$ of a simplicial complex $K$, and its
rank $\beta_k = \mathrm{rank}(H_k)$, the \emph{$k$-th Betti number}.

\paragraph{$H_0$: connected components.}
$\beta_0$ counts the number of connected components in the complex.
A complex with five isolated vertices has $\beta_0 = 5$.
Adding an edge between two vertices merges two components into one, so
$\beta_0$ drops by one.
In terms of attention graphs: $\beta_0$ tells us how many disconnected
clusters of tokens exist at a given threshold.
A head that attends broadly (connecting most tokens) will have
$\beta_0 = 1$ at moderate $\varepsilon$;
a head that attends narrowly will keep $\beta_0$ high for longer.

\paragraph{$H_1$: loops.}
$\beta_1$ counts independent loops---closed paths of edges that do not
bound a filled triangle.
Think of a square made of four edges with no diagonal: it has one loop
($\beta_1 = 1$).
Add a diagonal to split it into two triangles and fill both: the loop
dies ($\beta_1 = 0$).
In attention graphs, a loop signals a cycle of mutual attention among a
group of tokens where no single token mediates all the connections.
\textcite{kushnareva2021} found that $H_1$ features were among the most
discriminative for distinguishing human-written from machine-generated
text---machine text produced systematically different loop structures.

\paragraph{$H_2$ and higher.}
$\beta_2$ counts enclosed voids (think of the hollow interior of a
tetrahedron whose four triangular faces are present but whose
interior is not filled).
For the token counts typical of our attention graphs (sentences of
20--128 tokens), $H_2$ features are sparse and noisy.
We compute them when Ripser provides them, but they carry little signal
in practice.
Throughout this project, our workhorse features come from $H_0$ and
$H_1$.


\section{Persistence: birth, death, and the diagram}
\label{sec:persistence}

Now we combine the filtration with homology.
As $\varepsilon$ sweeps from $0$ upward, topological features appear and
disappear.
A connected component is \emph{born} at the smallest $\varepsilon$ where
its vertex first enters the complex, and it \emph{dies} when it merges
with an older component (by convention, the younger component dies---the
\emph{elder rule}).
A loop is born when its closing edge appears, and it dies when a
triangle fills it in.

Each feature yields a \emph{persistence pair} $(b_i, d_i)$, where $b_i$
is the birth value and $d_i$ is the death value.
The quantity $d_i - b_i$ is the feature's \emph{persistence}---how long
it survives across the filtration.
High-persistence features reflect robust structure in the data;
low-persistence features are typically noise or artifacts of the
particular sample \parencite{ghrist2008}.

Two equivalent visualizations represent the collection of persistence
pairs:

\begin{enumerate}
  \item \textbf{Persistence diagram.}
        Plot each pair $(b_i, d_i)$ as a point in the plane, with
        birth on the horizontal axis and death on the vertical axis.
        Every point lies above the diagonal $b = d$ (a feature always
        dies after it is born).
        Points far from the diagonal are high-persistence; points
        near it are noise.
  \item \textbf{Barcode.}
        Draw each pair as a horizontal bar from $b_i$ to $d_i$.
        Long bars correspond to persistent features; short bars to
        ephemeral ones.
        The barcode is the multiset view of the same information.
\end{enumerate}

% persistence-diagram.tex — example persistence diagram (birth vs death)
% Include via % persistence-diagram.tex — example persistence diagram (birth vs death)
% Include via % persistence-diagram.tex — example persistence diagram (birth vs death)
% Include via \input{figures/persistence-diagram} from overview.tex

\begin{figure}[H]
\centering
\begin{tikzpicture}[
    >=Stealth,
    h0point/.style={circle, fill=blue!70!black, inner sep=2.5pt},
    h1point/.style={diamond, fill=red!70!black, inner sep=2.5pt},
    annotation/.style={font=\small, anchor=west},
]

% Axes
\draw[->, thick] (-0.3, 0) -- (5.5, 0)
    node[anchor=north, font=\small] {birth};
\draw[->, thick] (0, -0.3) -- (0, 5.5)
    node[anchor=east, font=\small] {death};

% Diagonal y = x
\draw[gray, dashed, thin] (0, 0) -- (5.2, 5.2);
\node[gray, font=\footnotesize, rotate=45, anchor=south] at (4.2, 4.2) {$y = x$};

% --- H_0 features (blue circles) ---
% Feature 1: born at 0, dies at 0.9 (short-lived component merge)
\node[h0point] (h0a) at (0, 0.9) {};

% Feature 2: born at 0, dies at 1.2
\node[h0point] (h0b) at (0, 1.2) {};

% Feature 3: born at 0, never dies (essential feature) — plot near top
\node[h0point] (h0c) at (0, 5.0) {};
\draw[blue!70!black, thick, dashed] (h0c) -- (0, 5.4);

% --- H_1 features (red diamonds) ---
% Feature 4: born at 0.9, dies at 1.2 (short-lived loop)
\node[h1point] (h1a) at (0.9, 1.2) {};

% Feature 5: born at 0.7, dies at 2.0 (longest-lived loop — the one we mark)
\node[h1point] (h1b) at (0.7, 2.0) {};

% --- Annotate longest-living feature ---
\draw[red!70!black, thin, ->] (1.8, 2.6) -- ($(h1b) + (0.15, 0.15)$);
\node[annotation, red!70!black] at (1.8, 2.6) {longest-lived $H_1$ feature};

% --- Annotate essential component ---
\draw[blue!70!black, thin, ->] (1.2, 4.8) -- ($(h0c) + (0.15, 0.0)$);
\node[annotation, blue!70!black] at (1.2, 4.8) {essential component ($\infty$)};

% --- Persistence as vertical distance from diagonal ---
% Show for the longest-lived H_1 feature
\draw[red!50!black, thin, dotted] (0.7, 0.7) -- (h1b);
\node[font=\scriptsize, red!50!black, anchor=west] at (0.85, 1.35) {persistence};

% Axis tick marks
\foreach \x in {1, 2, 3, 4, 5} {
    \draw (\x, 0.07) -- (\x, -0.07) node[below, font=\scriptsize] {\x};
    \draw (0.07, \x) -- (-0.07, \x) node[left, font=\scriptsize] {\x};
}

% Legend
\node[h0point, label={[font=\footnotesize]right:$H_0$ (components)}]
    at (3.5, 4.8) {};
\node[h1point, label={[font=\footnotesize]right:$H_1$ (loops)}]
    at (3.5, 4.3) {};

\end{tikzpicture}
\caption{%
    A persistence diagram plots each topological feature as a point at
    coordinates (birth, death). Features far from the diagonal $y = x$
    persist across a wide range of $\varepsilon$ and represent robust
    structure; features near the diagonal are short-lived noise. $H_0$
    features (blue circles) track connected components; $H_1$ features (red
    diamonds) track loops. The dashed extension marks an essential component
    that never dies.%
}
\label{fig:persistence-diagram}
\end{figure}
 from overview.tex

\begin{figure}[H]
\centering
\begin{tikzpicture}[
    >=Stealth,
    h0point/.style={circle, fill=blue!70!black, inner sep=2.5pt},
    h1point/.style={diamond, fill=red!70!black, inner sep=2.5pt},
    annotation/.style={font=\small, anchor=west},
]

% Axes
\draw[->, thick] (-0.3, 0) -- (5.5, 0)
    node[anchor=north, font=\small] {birth};
\draw[->, thick] (0, -0.3) -- (0, 5.5)
    node[anchor=east, font=\small] {death};

% Diagonal y = x
\draw[gray, dashed, thin] (0, 0) -- (5.2, 5.2);
\node[gray, font=\footnotesize, rotate=45, anchor=south] at (4.2, 4.2) {$y = x$};

% --- H_0 features (blue circles) ---
% Feature 1: born at 0, dies at 0.9 (short-lived component merge)
\node[h0point] (h0a) at (0, 0.9) {};

% Feature 2: born at 0, dies at 1.2
\node[h0point] (h0b) at (0, 1.2) {};

% Feature 3: born at 0, never dies (essential feature) — plot near top
\node[h0point] (h0c) at (0, 5.0) {};
\draw[blue!70!black, thick, dashed] (h0c) -- (0, 5.4);

% --- H_1 features (red diamonds) ---
% Feature 4: born at 0.9, dies at 1.2 (short-lived loop)
\node[h1point] (h1a) at (0.9, 1.2) {};

% Feature 5: born at 0.7, dies at 2.0 (longest-lived loop — the one we mark)
\node[h1point] (h1b) at (0.7, 2.0) {};

% --- Annotate longest-living feature ---
\draw[red!70!black, thin, ->] (1.8, 2.6) -- ($(h1b) + (0.15, 0.15)$);
\node[annotation, red!70!black] at (1.8, 2.6) {longest-lived $H_1$ feature};

% --- Annotate essential component ---
\draw[blue!70!black, thin, ->] (1.2, 4.8) -- ($(h0c) + (0.15, 0.0)$);
\node[annotation, blue!70!black] at (1.2, 4.8) {essential component ($\infty$)};

% --- Persistence as vertical distance from diagonal ---
% Show for the longest-lived H_1 feature
\draw[red!50!black, thin, dotted] (0.7, 0.7) -- (h1b);
\node[font=\scriptsize, red!50!black, anchor=west] at (0.85, 1.35) {persistence};

% Axis tick marks
\foreach \x in {1, 2, 3, 4, 5} {
    \draw (\x, 0.07) -- (\x, -0.07) node[below, font=\scriptsize] {\x};
    \draw (0.07, \x) -- (-0.07, \x) node[left, font=\scriptsize] {\x};
}

% Legend
\node[h0point, label={[font=\footnotesize]right:$H_0$ (components)}]
    at (3.5, 4.8) {};
\node[h1point, label={[font=\footnotesize]right:$H_1$ (loops)}]
    at (3.5, 4.3) {};

\end{tikzpicture}
\caption{%
    A persistence diagram plots each topological feature as a point at
    coordinates (birth, death). Features far from the diagonal $y = x$
    persist across a wide range of $\varepsilon$ and represent robust
    structure; features near the diagonal are short-lived noise. $H_0$
    features (blue circles) track connected components; $H_1$ features (red
    diamonds) track loops. The dashed extension marks an essential component
    that never dies.%
}
\label{fig:persistence-diagram}
\end{figure}
 from overview.tex

\begin{figure}[H]
\centering
\begin{tikzpicture}[
    >=Stealth,
    h0point/.style={circle, fill=blue!70!black, inner sep=2.5pt},
    h1point/.style={diamond, fill=red!70!black, inner sep=2.5pt},
    annotation/.style={font=\small, anchor=west},
]

% Axes
\draw[->, thick] (-0.3, 0) -- (5.5, 0)
    node[anchor=north, font=\small] {birth};
\draw[->, thick] (0, -0.3) -- (0, 5.5)
    node[anchor=east, font=\small] {death};

% Diagonal y = x
\draw[gray, dashed, thin] (0, 0) -- (5.2, 5.2);
\node[gray, font=\footnotesize, rotate=45, anchor=south] at (4.2, 4.2) {$y = x$};

% --- H_0 features (blue circles) ---
% Feature 1: born at 0, dies at 0.9 (short-lived component merge)
\node[h0point] (h0a) at (0, 0.9) {};

% Feature 2: born at 0, dies at 1.2
\node[h0point] (h0b) at (0, 1.2) {};

% Feature 3: born at 0, never dies (essential feature) — plot near top
\node[h0point] (h0c) at (0, 5.0) {};
\draw[blue!70!black, thick, dashed] (h0c) -- (0, 5.4);

% --- H_1 features (red diamonds) ---
% Feature 4: born at 0.9, dies at 1.2 (short-lived loop)
\node[h1point] (h1a) at (0.9, 1.2) {};

% Feature 5: born at 0.7, dies at 2.0 (longest-lived loop — the one we mark)
\node[h1point] (h1b) at (0.7, 2.0) {};

% --- Annotate longest-living feature ---
\draw[red!70!black, thin, ->] (1.8, 2.6) -- ($(h1b) + (0.15, 0.15)$);
\node[annotation, red!70!black] at (1.8, 2.6) {longest-lived $H_1$ feature};

% --- Annotate essential component ---
\draw[blue!70!black, thin, ->] (1.2, 4.8) -- ($(h0c) + (0.15, 0.0)$);
\node[annotation, blue!70!black] at (1.2, 4.8) {essential component ($\infty$)};

% --- Persistence as vertical distance from diagonal ---
% Show for the longest-lived H_1 feature
\draw[red!50!black, thin, dotted] (0.7, 0.7) -- (h1b);
\node[font=\scriptsize, red!50!black, anchor=west] at (0.85, 1.35) {persistence};

% Axis tick marks
\foreach \x in {1, 2, 3, 4, 5} {
    \draw (\x, 0.07) -- (\x, -0.07) node[below, font=\scriptsize] {\x};
    \draw (0.07, \x) -- (-0.07, \x) node[left, font=\scriptsize] {\x};
}

% Legend
\node[h0point, label={[font=\footnotesize]right:$H_0$ (components)}]
    at (3.5, 4.8) {};
\node[h1point, label={[font=\footnotesize]right:$H_1$ (loops)}]
    at (3.5, 4.3) {};

\end{tikzpicture}
\caption{%
    A persistence diagram plots each topological feature as a point at
    coordinates (birth, death). Features far from the diagonal $y = x$
    persist across a wide range of $\varepsilon$ and represent robust
    structure; features near the diagonal are short-lived noise. $H_0$
    features (blue circles) track connected components; $H_1$ features (red
    diamonds) track loops. The dashed extension marks an essential component
    that never dies.%
}
\label{fig:persistence-diagram}
\end{figure}


\autoref{fig:persistence-diagram} shows both representations for a small
filtration.
In our pipeline, we compute a persistence diagram for each attention
head, separately for $H_0$ and $H_1$.
The collection of diagrams across all 144 attention heads of mBERT
(12 layers $\times$ 12 heads) forms the topological fingerprint of a
given input.


\section{The stability theorem}
\label{sec:stability}

Persistent homology would be a curiosity rather than a tool if its
output were fragile---if a tiny change in the input data could
rearrange the persistence diagram entirely.
The \emph{stability theorem} of \textcite{cohensteiner2007} guarantees
that this does not happen.

The theorem says: if two functions (in our case, two distance matrices
derived from attention weights) differ by at most $\delta$ in the
supremum norm, then their persistence diagrams differ by at most
$\delta$ in the \emph{bottleneck distance} $d_B$.
Formally, for tame functions $f$ and $g$,
\[
  d_B\bigl(\mathrm{Dgm}(f),\, \mathrm{Dgm}(g)\bigr) \;\leq\; \|f - g\|_\infty.
\]

The bottleneck distance $d_B$ between two diagrams is defined as the
infimum, over all bijections $\gamma$ between the two multisets of
points (with points on the diagonal available as padding), of the
largest distance between matched pairs:
\[
  d_B(D, D') = \inf_\gamma \sup_i \| p_i - \gamma(p_i) \|_\infty.
\]

Why does stability matter for us?
Attention weights vary slightly across runs of the same input
(dropout during inference, floating-point nondeterminism on GPU).
Stability guarantees that these small perturbations produce
small perturbations of the diagram.
The topological features we extract are not artifacts of numerical
noise---they reflect genuine structure in the attention pattern.
This is the property that makes persistent homology a reliable feature
extractor rather than an over-sensitive novelty.


\section{Comparing persistence diagrams}
\label{sec:comparing-diagrams}

The bottleneck distance $d_B$ described above is one way to compare two
persistence diagrams; it measures the worst-case mismatch.
A complementary metric is the \emph{$p$-Wasserstein distance}:
\[
  W_p(D, D') = \left(\inf_\gamma \sum_i \| p_i - \gamma(p_i)
  \|_\infty^p \right)^{1/p},
\]
which sums over all matched pairs rather than taking the supremum.
For $p = 2$, $W_2$ is particularly common.
The bottleneck distance is sensitive to the single most different
feature; the Wasserstein distance is sensitive to the aggregate of all
differences.
Both have stability guarantees \parencite{cohensteiner2007}.

These distances operate directly on diagrams, which is elegant but
inconvenient for machine learning: most classifiers expect
fixed-dimensional vector inputs, and persistence diagrams have variable
numbers of points.
Two vectorization strategies bridge the gap:

\paragraph{Persistence landscapes.}
\textcite{bubenik2015} proposed mapping each diagram to a sequence of
piecewise-linear functions (``landscapes'') that live in a Banach space.
Landscapes support averaging, statistical tests, and direct use as
feature vectors.

\paragraph{Persistence images.}
\textcite{adams2017} instead place a Gaussian kernel at each point of
the diagram, weight by a function of persistence, and discretize the
result onto a grid.
The output is a fixed-size matrix---an image---that can be fed to any
classifier.

\textcite{kushnareva2021} took a different and more direct approach:
rather than vectorizing diagrams via landscapes or images, they
extracted hand-crafted summary statistics from the barcodes---mean bar
length, variance, entropy, counts above thresholds---and concatenated
these across attention heads.
This trades representational richness for interpretability and
computational simplicity, and it worked well enough to achieve
state-of-the-art results on artificial text detection.
Our adaptation inherits this choice (\autoref{ch:pipeline} walks through
the specific features).


\section{What features get extracted}
\label{sec:features-conceptual}

Before we defer the full feature specification to \autoref{ch:pipeline},
a conceptual preview will help orient the reader.
\textcite{kushnareva2021} extract three families of topological features
from each attention head:

\paragraph{Barcode statistics.}
From the $H_0$ and $H_1$ barcodes, they compute: the number of bars,
the sum, mean, and variance of bar lengths, the birth and death times of
the longest bar, the number of bars born before or after a given
threshold, and the entropy of the barcode (treating normalized bar
lengths as a probability distribution).
These are scalar summaries of the persistence diagram.

\paragraph{Threshold features.}
At each of several fixed thresholds $\varepsilon_1 < \cdots <
\varepsilon_k$, they record the number of edges, the number of connected
components ($\beta_0$), and the number of strongly connected components
and simple cycles in the directed version of the graph.
These capture the graph's structure at specific scales rather than
summarizing across the whole filtration.

\paragraph{Pattern-distance features.}
\textcite{clark2019} identified recurring structural motifs in BERT
attention: attention to the previous token, to the next token, to
\texttt{[CLS]}, to \texttt{[SEP]}, and to punctuation.
Kushnareva formalize each pattern as a template graph and measure the
(normalized Frobenius) distance from the actual attention graph to each
template at every threshold.
These features quantify how closely a head's attention conforms to known
patterns.

All three families are concatenated across the 144 attention heads of
BERT to form a single high-dimensional feature vector per input text.
For our cross-linguistic comparison, the question becomes: do these
topological fingerprints cluster by language, by semantic domain, or by
both?


\section{Computation: Ripser and giotto-tda}
\label{sec:computation}

Computing persistent homology reduces to a matrix reduction problem
over $\mathbb{Z}/2\mathbb{Z}$---the boundary matrices of the simplicial
complex are reduced to a normal form that directly yields the
persistence pairs \parencite{zomorodian2005}.
The naive approach is cubic in the number of simplices, which is already
exponential in the number of vertices.
For the Vietoris--Rips complex on $n$ points, the complex can have up to
$2^n$ simplices, so efficient algorithms are essential.

\textcite{bauer2021} developed \emph{Ripser}, an algorithm and
implementation that exploits three key optimizations:
the \emph{clearing} optimization (many columns reduce to zero and can be
skipped), \emph{cohomology} (working with coboundary matrices, which are
sparser than boundary matrices in the Rips filtration), and an
\emph{implicit} representation of the Vietoris--Rips complex that avoids
enumerating all simplices.
Ripser is the standard tool for Vietoris--Rips persistent homology
and is typically an order of magnitude faster than alternatives.

In our pipeline, we access Ripser through two Python interfaces.
The \texttt{ripser} Python package (version 0.6.14 in our environment)
wraps the C++ core for CPU computation.
\texttt{giotto-tda} \parencite{tauzin2021} provides a scikit-learn-compatible
API with transformers for persistence diagrams, Betti curves, and
landscapes.
For GPU-accelerated computation, we use \texttt{ripserplusplus}---a CUDA
port that is substantially faster on large inputs but requires careful
version pinning.
(The PyPI release does not compile against CUDA~12+; we pin to a
specific commit.
\autoref{ch:tooling} documents the details.)

The computational cost is manageable for our setting.
Each attention head produces a matrix of at most $128 \times 128$
(the maximum sequence length after tokenization), and Ripser handles
matrices of this size in milliseconds.
The full pipeline across all 144 heads, three languages, and three
semantic domains remains well within the capacity of a single GPU
workstation.
